% Source ownership:
% - Itxaso: scaling of replicas + observables
% - Arthur/Manel: broader discussion and combined-program logic

\section{Scaling Analysis and Discussion}
\label{sec:scaling}

\subsection{Replica Parallelism: Speedup vs. System Size and Run Length}

The independent-replica MPI strategy was benchmarked by varying both the
polymer size and the production run length. In this approach, each MPI
worker executes a different production replica, so the total wall time
depends primarily on how much useful work can be distributed across
workers before communication and orchestration overhead begin to matter.

The first scaling study varies the number of carbons while keeping the
simulation length fixed. As the chain becomes larger, each replica
contains more internal work, especially in the energy calculations and
observable evaluations, so the MPI orchestration overhead becomes less
significant relative to the total runtime. For this reason, larger
systems are expected to display better parallel efficiency than very
small ones.

The second study varies the number of Monte Carlo steps for a fixed
system size. Short runs do not amortize the startup, scheduling, and I/O
costs very well, whereas longer runs allow each worker to spend a much
greater fraction of time performing actual MC sampling. Consequently,
speedup generally improves as the number of production steps increases.

In both cases, the speedup is not expected to be perfectly linear. The
master node still performs administrative tasks such as initial file
reading, dispatching jobs, receiving completion messages, and managing
the production queue. This means that part of the total hardware budget
is always reserved for orchestration rather than direct numerical work.

\begin{figure}[h!]
  \centering
  \includegraphics[width=\textwidth]{speedup_vs_natoms.pdf}
  \caption{Speedup of independent-replica parallelization: speedup as a function of chain length
           (\(N_C \in \{20,50,100,200,500\}\), \(n_{\text{steps}} = 10^6\)).
           All runs use 3 MPI ranks. The dashed line indicates ideal
           linear scaling (\(S = 3\)).}
\end{figure}

\begin{figure}[h!]
  \centering
  \includegraphics[width=\textwidth]{speedup_vs_nsteps.pdf}
  \caption{Speedup of independent-replica parallelization: speedup as a function of run length
           (\(n_{\text{steps}} \in \{10^5, 10^6, 10^7\}\), \(N_C = 100\)).
           All runs use 3 MPI ranks. The dashed line indicates ideal
           linear scaling (\(S = 3\)).}
\end{figure}

\subsection{Observable Post-Processing: Strong Scaling}

A second scaling study was carried out for the observable
post-processing kernel, implemented in
\texttt{main\_parallel\_observables}. Unlike the production simulations,
this stage is a strong-scaling problem: the total dataset is fixed, and
the number of MPI processes is increased to measure how efficiently the
same workload can be completed.

Each process reads a subset of trajectory files, computes the radius of
gyration, end-to-end distance, and torsional statistics, and
accumulates local partial sums. At the end of the run, the global
results are assembled through MPI reductions on rank 0, which writes the
final averaged observables.

This workload is well suited to MPI because each trajectory file can be
processed independently for most of the runtime. Communication is only
required at the end, when the local statistics are combined into a
single global dataset. As a result, the observable analysis behaves much
closer to an embarrassingly parallel workload than the simulation phase
itself.

The scaling is therefore expected to be strong for moderate process
counts, with efficiency gradually decreasing once the number of workers
becomes large compared to the number of files. At that point, the amount
of useful work per process shrinks and the fixed costs of process launch,
file I/O, and final reductions become more visible.

\begin{figure}[h!]
  \centering
  \includegraphics[width=\textwidth]{benchmark_vs_np.pdf}
  \caption{Strong scaling of the observable post-processing kernel.
           MPI wall-clock time (solid circles) and shell-measured
           elapsed time (dashed squares) as a function of the number of
           MPI processes \(n_p\).}
\end{figure}

\begin{figure}[h!]
  \centering
  \includegraphics[width=\textwidth]{observables_vs_np.pdf}
  \caption{Mean radius of gyration \(R_g\) and mean end-to-end distance
           \(R_{ee}\) as a function of the number of MPI processes \(n_p\).
           The overlap of all curves confirms that the observable averages are
           preserved exactly across the tested parallel decompositions.}
\end{figure}

\subsection{Discussion}

One of the key observations from our performance analysis is that
parallelization rarely translates into perfectly linear speedup. In an
ideal case, doubling the number of cores would halve the runtime, but in
practice communication overhead, process coordination, file handling,
and hardware contention all reduce the efficiency of additional workers.

This effect is visible in both the independent-replica and
post-processing benchmarks. For the MPI production workflow, the master
must coordinate task dispatch and receive status messages from workers.
For the observables workflow, the amount of useful work per process
eventually becomes too small compared to the fixed costs of launching
processes and reducing the final results.

Resource utilization is therefore just as important as raw parallelism.
The orchestrator-worker design helps maintain high CPU occupancy by
ensuring that free workers are immediately assigned new tasks from the
queue. At the same time, the staged execution model avoids wasting
resources on production jobs before equilibration has completed.

Combining MPI and OpenMP also requires careful control to avoid
oversubscription. If too many MPI ranks each spawn too many OpenMP
threads, the runtime can degrade due to thread collisions and operating
system scheduling overhead. This is why the Makefile parameters
\texttt{NP} and \texttt{OMPTHREADS} are treated as a coordinated
resource-management mechanism rather than as independent tuning knobs.

Finally, the scaling trends also reflect the nonlinearity of Monte Carlo
workloads themselves. Some replicas converge faster than others, some
trajectories generate more favourable moved/fixed partitions inside
\texttt{delta\_energy}, and some process counts match the available job
granularity better than others. For this reason, the best-performing
configuration is not always the one with the largest number of cores,
but the one that best matches the size and structure of the workload.

\subsection{Combined Parallel Workflow}

The final program combines several levels of parallelism into a single
coordinated workflow. At the top level, MPI is used to distribute work
across processes in an orchestrator-worker star topology. Depending on
the phase of execution, workers either participate in collaborative
equilibration or execute independent production replicas drawn from a
global queue.

Within each worker, OpenMP is used to accelerate the most expensive
energy kernel, the nested Lennard-Jones loop in
\texttt{delta\_energy}. This creates a hybrid structure in which MPI
handles macro-level task distribution while OpenMP performs
micro-level acceleration inside each Monte Carlo step.


A key design choice is that the simulation is separated into phases so
that resources can be reused dynamically. If an equilibrated geometry
already exists, the master can bypass the equilibration stage for that
configuration and immediately unlock the corresponding production jobs.
Meanwhile, workers that finish equilibration for other conformations are
reassigned to the shared production queue, ensuring that CPU resources
are not left idle.

This combined architecture does not simply attempt to maximize the
number of active cores at all times. Instead, it tries to maximize the
amount of useful work completed per unit time while respecting the
constraints of communication overhead, thread contention, and workload
granularity. In this sense, the program acts as a coordinated HPC
ecosystem rather than as a single parallel loop.
